% TOP_PROBLEM: {"id": 2747, "title": "Kirby Problem 1.88", "category_id": 7, "category": {"id": 7, "name": "topology", "display_name": "Topology", "description": "Properties preserved under continuous deformations.", "slug": "topology", "order_index": 7, "created_at": "2026-07-31T15:26:25.671Z"}, "status": "open", "problem_number": "KP-1.88", "set_id": 11, "statusQualification": "Restores the imaginary quadratic field and counts arithmetic links by their complements, as in the source's lattice formulation."}
% FIRST_POSTED: 2026-10-01
% ENTRY_KIND: statement_only
% UnsolvedMath ID 2747; source catalogue code KP-1.88
% !TeX program = lualatex
% Definition and sources only; this file is not an LLM solution attempt.
\documentclass[11pt,a4paper]{article}
\usepackage[margin=25mm]{geometry}
\usepackage{fontspec}
\setmainfont{lmroman10-regular.otf}[BoldFont=lmroman10-bold.otf,
 ItalicFont=lmroman10-italic.otf,BoldItalicFont=lmroman10-bolditalic.otf]
\usepackage{amsmath,amssymb,mathtools}
\usepackage{unicode-math}
\setmathfont{latinmodern-math.otf}
\AtBeginDocument{\let\setminus\smallsetminus}
\usepackage{xurl,hyperref}
\hypersetup{colorlinks=true,urlcolor=blue}
\setcounter{secnumdepth}{0}
\setlength{\parindent}{0pt}
\setlength{\parskip}{5pt}
\setlength{\emergencystretch}{3em}
\begin{document}
\section*{Kirby Problem 1.88}\label{2747}
{\small UnsolvedMath ID 2747 · KP-1.88 · Topology · Kirby's Problems in Low-Dimensional Topology}
\subsection{Definitions and mathematical statement}
Let $d$ be a positive squarefree integer and let $\mathcal O_d$ be the ring of integers of the imaginary quadratic field $\mathbb Q(\sqrt{-d})$. Its Bianchi group $\operatorname{PSL}_2(\mathcal O_d)$ acts by orientation-preserving isometries on hyperbolic three-space. For a nonzero ideal $I\subset\mathcal O_d$, define the principal congruence subgroup
\[
\Gamma(I)=\ker\bigl(\operatorname{PSL}_2(\mathcal O_d)
\longrightarrow\operatorname{PSL}_2(\mathcal O_d/I)\bigr).
\]
A subgroup $\Gamma\le\operatorname{PSL}_2(\mathcal O_d)$ is congruence when it contains some $\Gamma(I)$; equality is not required.

A congruence arithmetic link complement is a manifold $S^3\setminus L$ homeomorphic to $\mathbb H^3/\Gamma$ for such a torsion-free congruence subgroup, allowing conjugation of the lattice in $\operatorname{PSL}_2(\mathbb C)$. Here $L$ is a finite smooth link and the quotient has its complete finite-volume metric. Ask whether there are infinitely many such link complements, up to homeomorphism, as $d$, $I$, and $\Gamma$ vary.

This is the geometric meaning of infinitude of congruence arithmetic links in the source: it counts distinct complements, rather than repeated diagrams or different choices of congruence data for one manifold. A principal congruence example belongs to the class, but restricting to principal subgroups would change the problem. Likewise, commensurability with a Bianchi group without the congruence containment is insufficient.
\subsection{Short English statement}
Are there infinitely many distinct link complements in the sphere whose hyperbolic groups are congruence subgroups of Bianchi groups?
\subsection{Sources}
\begin{itemize}
\item[S1] ULAM.AI and the UnsolvedMath Contributors, supplied problems.json, record 2747 (KP-1.88), Kirby's Problems in Low-Dimensional Topology. Catalogue identity and source transcription; CC BY 4.0.
\url{https://huggingface.co/datasets/ulamai/UnsolvedMath}.
\item[S2] K3: A New Problem List in Low-Dimensional Topology, Problem 1.88. Expanded formulation of the supplied catalogue statement.
\url{https://math.berkeley.edu/sites/default/files/surv-295-ruberman-watermarked-author-pdf.pdf}.
\end{itemize}
\end{document}
